% \Act*{4}

% \begin{figure}[t]
	% \centering
	% \includegraphics[width=\textwidth]{act4head}
% \end{figure}

\Scene{1}[The Forest of Arden]

\enter{\textsc{Rosalind} as \textsc{Ganymede,} \textsc{Celia} as \textsc{Aliena,} and \textsc{Jaques.}}


\begin{prose_speech}[Jaques]
I prithee, pretty youth, let me be better acquainted with thee.
\end{prose_speech}

\begin{prose_speech}[Rosalind as Ganymede]
They say you are a melancholy fellow.
\end{prose_speech}

\begin{prose_speech}[Jaques]
I am so. I do love it better than laughing.
\end{prose_speech}

\begin{prose_speech}[Rosalind as Ganymede]
Those that are in extremity of either are abominable fellows and betray themselves to every modern censure worse than drunkards.
\end{prose_speech}

\begin{prose_speech}[Jaques]
Why, 'tis good to be sad and say nothing.
\end{prose_speech}

\begin{prose_speech}[Rosalind as Ganymede]
Why then, 'tis good to be a post.
\end{prose_speech}

\begin{prose_speech}[Jaques]
I have neither the scholar's melancholy, which is emulation; nor the musician's, which is fantastical; nor the courtier's, which is proud; nor the soldier's, which is ambitious; nor the lawyer's, which is politic; nor the lady's, which is nice; nor the lover's, which is all these; but it is a melancholy of mine own, compounded of many simples, extracted from many objects, and indeed the sundry contemplation of my travels, in which my often rumination wraps me in a most humorous sadness.
\end{prose_speech}

\begin{prose_speech}[Rosalind as Ganymede]
A traveller. By my faith, you have great reason to be sad. I fear you have sold your own lands to see other men's. Then to have seen much and to have nothing is to have rich eyes and poor hands.
\end{prose_speech}

\begin{prose_speech}[Jaques]
Yes, I have gained my experience.
\end{prose_speech}

\begin{prose_speech}[Rosalind as Ganymede]
And your experience makes you sad. I had rather have a fool to make me merry than experience to make me sad—and to travel for it too.
\end{prose_speech}

\enter{\textsc{Orlando}}


\begin{prose_speech}[Orlando]
Good day and happiness, dear Rosalind.
\end{prose_speech}

\begin{prose_speech}[Jaques]
Nay then, God be wi' you, an you talk in blank verse.
\end{prose_speech}

\begin{prose_speech}[Rosalind as Ganymede]
Farewell, Monsieur Traveller. Look you lisp and wear strange suits, disable all the benefits of your own country, be out of love with your nativity, and almost chide God for making you that countenance you are, or I will scarce think you have swam in a gondola.
\exit{\textsc{Jaques}}
Why, how now, Orlando, where have you been all this while? You a lover? An you serve me such another trick, never come in my sight more.
\end{prose_speech}

\begin{prose_speech}[Orlando]
My fair Rosalind, I come within an hour of my promise.
\end{prose_speech}

\begin{prose_speech}[Rosalind as Ganymede]
Break an hour's promise in love? He that will divide a minute into a thousand parts and break but a part of the thousand part of a minute in the affairs of love, it may be said of him that Cupid hath clapped him o' th' shoulder, but I'll warrant him heart-whole.
\end{prose_speech}

\begin{prose_speech}[Orlando]
Pardon me, dear Rosalind.
\end{prose_speech}

\begin{prose_speech}[Rosalind as Ganymede]
Nay, an you be so tardy, come no more in my sight. I had as lief be wooed of a snail.
\end{prose_speech}

\begin{prose_speech}[Orlando]
Of a snail?
\end{prose_speech}

\begin{prose_speech}[Rosalind as Ganymede]
Ay, of a snail, for though he comes slowly, he carries his house on his head—a better jointure, I think, than you make a woman. Besides, he brings his destiny with him.
\end{prose_speech}

\begin{prose_speech}[Orlando]
What's that?
\end{prose_speech}

\begin{prose_speech}[Rosalind as Ganymede]
Why, horns, which such as you are fain to be beholding to your wives for. But he comes armed in his fortune and prevents the slander of his wife.
\end{prose_speech}

\begin{prose_speech}[Orlando]
Virtue is no hornmaker, and my Rosalind is virtuous.
\end{prose_speech}

\begin{prose_speech}[Rosalind as Ganymede]
And I am your Rosalind.
\end{prose_speech}

\begin{prose_speech}[Celia as Aliena]
It pleases him to call you so, but he hath a Rosalind of a better leer than you.
\end{prose_speech}

\begin{prose_speech}[Rosalind as Ganymede]
	\stage{To \textsc{Orlando}}  Come, woo me, woo me, for now I am in a holiday humour, and like enough to consent. What would you say to me now an I were your very, very Rosalind?
\end{prose_speech}

\begin{prose_speech}[Orlando]
I would kiss before I spoke.
\end{prose_speech}

\begin{prose_speech}[Rosalind as Ganymede]
Nay, you were better speak first, and when you were gravelled for lack of matter, you might take occasion to kiss. Very good orators, when they are out, they will spit; and for lovers lacking—God warn us—matter, the cleanliest shift is to kiss.
\end{prose_speech}

\begin{prose_speech}[Orlando]
How if the kiss be denied?
\end{prose_speech}

\begin{prose_speech}[Rosalind as Ganymede]
Then she puts you to entreaty, and there begins new matter.
\end{prose_speech}

\begin{prose_speech}[Orlando]
Who could be out, being before his beloved mistress?
\end{prose_speech}

\begin{prose_speech}[Rosalind as Ganymede]
Marry, that should you if I were your mistress, or I should think my honesty ranker than my wit.
\end{prose_speech}

\begin{prose_speech}[Orlando]
What, of my suit?
\end{prose_speech}

\begin{prose_speech}[Rosalind as Ganymede]
Not out of your apparel, and yet out of your suit. Am not I your Rosalind?
\end{prose_speech}

\begin{prose_speech}[Orlando]
I take some joy to say you are because I would be talking of her.
\end{prose_speech}

\begin{prose_speech}[Rosalind as Ganymede]
Well, in her person I say I will not have you.
\end{prose_speech}

\begin{prose_speech}[Orlando]
Then, in mine own person I die.
\end{prose_speech}

\begin{prose_speech}[Rosalind as Ganymede]
No, faith, die by attorney. The poor world is almost six thousand years old, and in all this time there was not any man died in his own person, videlicet, in a love cause. Troilus had his brains dashed out with a Grecian club, yet he did what he could to die before, and he is one of the patterns of love. Leander, he would have lived many a fair year though Hero had turned nun, if it had not been for a hot midsummer night, for, good youth, he went but forth to wash him in the Hellespont and, being taken with the cramp, was drowned; and the foolish chroniclers of that age found it was Hero of Sestos. But these are all lies. Men have died from time to time and worms have eaten them, but not for love.
\end{prose_speech}

\begin{prose_speech}[Orlando]
I would not have my right Rosalind of this mind, for I protest her frown might kill me.
\end{prose_speech}

\begin{prose_speech}[Rosalind as Ganymede]
By this hand, it will not kill a fly. But come; now I will be your Rosalind in a more coming-on disposition, and ask me what you will, I will grant it.
\end{prose_speech}

\begin{colorbigpic}
	[\basicscale]
	{4i_comingon}
	[I will be your Rosalind in a more coming-on disposition]
\end{colorbigpic}

\begin{prose_speech}[Orlando]
Then love me, Rosalind.
\end{prose_speech}

\begin{prose_speech}[Rosalind as Ganymede]
Yes, faith, will I, Fridays and Saturdays and all.
\end{prose_speech}

\begin{prose_speech}[Orlando]
And wilt thou have me?
\end{prose_speech}

\begin{prose_speech}[Rosalind as Ganymede]
Ay, and twenty such.
\end{prose_speech}

\begin{prose_speech}[Orlando]
What sayest thou?
\end{prose_speech}

\begin{prose_speech}[Rosalind as Ganymede]
Are you not good?
\end{prose_speech}

\begin{prose_speech}[Orlando]
I hope so.
\end{prose_speech}

\begin{prose_speech}[Rosalind as Ganymede]
Why then, can one desire too much of a good thing?—Come, sister, you shall be the priest and marry us.—Give me your hand, Orlando.—What do you say, sister?
\end{prose_speech}

\begin{prose_speech}[Orlando]
	\stage{To \textsc{Celia}}  Pray thee marry us.
\end{prose_speech}

\begin{prose_speech}[Celia as Aliena]
I cannot say the words.
\end{prose_speech}

\begin{prose_speech}[Rosalind as Ganymede]
You must begin <Will you, Orlando—>
\end{prose_speech}

\begin{prose_speech}[Celia as Aliena]
Go to.—Will you, Orlando, have to wife this Rosalind?
\end{prose_speech}

\begin{colorbigpic}
	[\basicscale]
	{4i_wife}
	[Will you, Orlando, have to wife this Rosalind?]
\end{colorbigpic}

\begin{prose_speech}[Orlando]
I will.
\end{prose_speech}

\begin{prose_speech}[Rosalind as Ganymede]
Ay, but when?
\end{prose_speech}

\begin{prose_speech}[Orlando]
Why now, as fast as she can marry us.
\end{prose_speech}

\begin{prose_speech}[Rosalind as Ganymede]
Then you must say <I take thee, Rosalind, for wife.>
\end{prose_speech}

\begin{prose_speech}[Orlando]
I take thee, Rosalind, for wife.
\end{prose_speech}

\begin{prose_speech}[Rosalind as Ganymede]
I might ask you for your commission, but I do take thee, Orlando, for my husband. There's a girl goes before the priest, and certainly a woman's thought runs before her actions.
\end{prose_speech}

\begin{prose_speech}[Orlando]
So do all thoughts. They are winged.
\end{prose_speech}

\begin{prose_speech}[Rosalind as Ganymede]
Now tell me how long you would have her after you have possessed her?
\end{prose_speech}

\begin{prose_speech}[Orlando]
Forever and a day.
\end{prose_speech}

\begin{prose_speech}[Rosalind as Ganymede]
Say <a day> without the <ever.> No, no, Orlando, men are April when they woo, December when they wed. Maids are May when they are maids, but the sky changes when they are wives. I will be more jealous of thee than a Barbary cock-pigeon over his hen, more clamorous than a parrot against rain, more newfangled than an ape, more giddy in my desires than a monkey. I will weep for nothing, like Diana in the fountain, and I will do that when you are disposed to be merry. I will laugh like a hyena, and that when thou art inclined to sleep.
\end{prose_speech}

\begin{prose_speech}[Orlando]
But will my Rosalind do so?
\end{prose_speech}

\begin{prose_speech}[Rosalind as Ganymede]
By my life, she will do as I do.
\end{prose_speech}

\begin{prose_speech}[Orlando]
O, but she is wise.
\end{prose_speech}

\begin{prose_speech}[Rosalind as Ganymede]
Or else she could not have the wit to do this. The wiser, the waywarder. Make the doors upon a woman's wit, and it will out at the casement. Shut that, and 'twill out at the keyhole. Stop that, 'twill fly with the smoke out at the chimney.
\end{prose_speech}

\begin{prose_speech}[Orlando]
A man that had a wife with such a wit, he might say <Wit, whither wilt?>
\end{prose_speech}

\begin{prose_speech}[Rosalind as Ganymede]
Nay, you might keep that check for it till you met your wife's wit going to your neighbour's bed.
\end{prose_speech}

\begin{prose_speech}[Orlando]
And what wit could wit have to excuse that?
\end{prose_speech}

\begin{prose_speech}[Rosalind as Ganymede]
Marry, to say she came to seek you there. You shall never take her without her answer unless you take her without her tongue. O, that woman that cannot make her fault her husband's occasion, let her never nurse her child herself, for she will breed it like a fool.
\end{prose_speech}

\begin{prose_speech}[Orlando]
For these two hours, Rosalind, I will leave thee.
\end{prose_speech}

\begin{prose_speech}[Rosalind as Ganymede]
Alas, dear love, I cannot lack thee two hours.
\end{prose_speech}

\begin{prose_speech}[Orlando]
I must attend the Duke at dinner. By two o'clock I will be with thee again.
\end{prose_speech}

\begin{prose_speech}[Rosalind as Ganymede]
Ay, go your ways, go your ways. I knew what you would prove. My friends told me as much, and I thought no less. That flattering tongue of yours won me. 'Tis but one cast away, and so, come, death. Two o'clock is your hour?
\end{prose_speech}

\begin{colorbigpic}
	[\basicscale]
	{4i_blind}
	[That blind rascally boy]
\end{colorbigpic}

\begin{prose_speech}[Orlando]
Ay, sweet Rosalind.
\end{prose_speech}

\begin{prose_speech}[Rosalind as Ganymede]
By my troth, and in good earnest, and so God mend me, and by all pretty oaths that are not dangerous, if you break one jot of your promise or come one minute behind your hour, I will think you the most pathetical break-promise, and the most hollow lover, and the most unworthy of her you call Rosalind that may be chosen out of the gross band of the unfaithful. Therefore beware my censure, and keep your promise.
\end{prose_speech}

\begin{prose_speech}[Orlando]
With no less religion than if thou wert indeed my Rosalind. So, adieu.
\end{prose_speech}



\begin{prose_speech}[Rosalind as Ganymede]
Well, time is the old justice that examines all such offenders, and let time try. Adieu.
\end{prose_speech}

\exit{\textsc{Orlando}}

\begin{prose_speech}[Celia]
You have simply misused our sex in your love-prate. We must have your doublet and hose plucked over your head and show the world what the bird hath done to her own nest.
\end{prose_speech}



\begin{prose_speech}[Rosalind]
O coz, coz, coz, my pretty little coz, that thou didst know how many fathom deep I am in love. But it cannot be sounded; my affection hath an unknown bottom, like the Bay of Portugal.
\end{prose_speech}

\begin{prose_speech}[Celia]
Or rather bottomless, that as fast as you pour affection in, it runs out.
\end{prose_speech}

\begin{prose_speech}[Rosalind]
No, that same wicked bastard of Venus, that was begot of thought, conceived of spleen, and born of madness, that blind rascally boy that abuses everyone's eyes because his own are out, let him be judge how deep I am in love. I'll tell thee, Aliena, I cannot be out of the sight of Orlando. I'll go find a shadow and sigh till he come.
\end{prose_speech}

\begin{prose_speech}[Celia]
And I'll sleep.
\end{prose_speech}

\exeunt{}
